\section*{Exercises about Span and Linear Independence}

\exercise{1}
\begin{xrcs}
  Find a list of four distinct vectors in $\myF^3$ whose span equals
  \begin{equation}
    \left \{  \left(x,y,z \right) \in \myF^3 \mid x + y + z = 0 \right \}.
  \end{equation}

  \begin{xsol}
    The plane $x+y+z=0$ is spanned by the two linearly independent vectors $(1,-1,0)$ and $(0,1,-1)$. To get four distinct vectors that span the same subspace, we can add two more distinct nonzero vectors from the plane, for instance:
    \begin{equation}
      (1,-1,0), \quad (0,1,-1), \quad (1,0,-1), \quad (2,-1,-1).
    \end{equation}
    Since each of these vectors satisfies $x+y+z=0$, their span is contained in the plane. Furthermore, since the first two vectors already span the entire plane, the span of all four vectors equals the desired set.
  \end{xsol}
\end{xrcs}

\exercise{3}
\begin{xrcs}
  Claim: If $v_1, v_2, v_3, v_4$ spans $V$, then the list $v_1 - v_2, v_2 - v_3, v_3 - v_4, v_4$ also spans $V$.

  \begin{xprf}
    Let $u \in V$. Since $v_1, v_2, v_3, v_4$ spans $V$, there exist scalars $b_1, b_2, b_3, b_4 \in \myF$ such that
    \[
      u = b_1 v_1 + b_2 v_2 + b_3 v_3 + b_4 v_4.
    \]
    We seek scalars $a_1, a_2, a_3, a_4 \in \myF$ such that
    \[
    \begin{aligned}
      u &= a_1 (v_1 - v_2) + a_2 (v_2 - v_3) + a_3 (v_3 - v_4) + a_4 v_4 \\
        &= a_1 v_1 + (a_2 - a_1) v_2 + (a_3 - a_2) v_3 + (a_4 - a_3) v_4.
    \end{aligned}
    \]
    Matching coefficients gives the system:
    \[
    \begin{aligned}
      a_1 &= b_1, \\
      a_2 - a_1 &= b_2 \iff a_2 = b_1 + b_2, \\
      a_3 - a_2 &= b_3 \iff a_3 = b_1 + b_2 + b_3, \\
      a_4 - a_3 &= b_4 \iff a_4 = b_1 + b_2 + b_3 + b_4.
    \end{aligned}
    \]
    Hence, any $u \in V$ can be written as
    \[
      u = b_1 (v_1 - v_2) + (b_1+b_2)(v_2 - v_3) + (b_1+b_2+b_3)(v_3 - v_4) + (b_1+b_2+b_3+b_4) v_4.
    \]
    Therefore, the list $v_1 - v_2, v_2 - v_3, v_3 - v_4, v_4$ also spans $V$.
  \end{xprf}
\end{xrcs}

\exercise{4}
\begin{xrcs}
  Claim: If $v_1, \ddd, v_m \in V$ and
  \[
    w_k := v_1 + \cdots + v_k \quad \forall k\in \{1, \ddd, m\},
  \]
  then $\myspan{v_1, \ddd, v_m} = \myspan{w_1, \ddd, w_m}$.

  \begin{xprf}
    Observe that
    \[
    \begin{aligned}
      w_1 &= v_1
      & \qquad v_1 &= w_1       \\
      w_2 &= v_1 + v_2
      & \qquad v_2 &= w_2 - w_1 \\
      w_3 &= v_1 + v_2 + v_3
      & \qquad v_3 &= w_3 - w_2 \\
      w_4 &= v_1 + v_2 + v_3 + v_4
      & \qquad v_4 &= w_4 - w_3 \\
      &\;\;\vdots
      &            &\;\;\vdots \\
      w_m &= v_1 + \cdots + v_m
      & \qquad v_m &= w_m - w_{m-1}.
    \end{aligned}
    \]

    \StepOne If $u \in \myspan{v_1, \ddd, v_m}$, there exist $a_1, \ddd, a_m \in \myF$ such that
    \[
    \begin{aligned}
      u &= a_1 v_1 + \cdots + a_m v_m \\
        &= a_1 w_1 + a_2 (w_2 - w_1) + a_3 (w_3 - w_2) + \cdots + a_m (w_m - w_{m-1}) \\
        &= (a_1 - a_2) w_1 + (a_2 - a_3) w_2 + \cdots + (a_{m-1} - a_m) w_{m-1} + a_m w_m.
    \end{aligned}
    \]
    Thus, $u \in \myspan{w_1, \ddd, w_m}$, and therefore,
    \begin{equation}
      \label{eq: fact1}
      \myspan{v_1, \ddd, v_m} \subseteq \myspan{w_1, \ddd, w_m}.
    \end{equation}

    \StepTwo If $u \in \myspan{w_1, \ddd, w_m}$, there exist $a_1, \ddd, a_m \in \myF$ such that
    \[
    \begin{aligned}
      u &= a_1 w_1 + \cdots + a_m w_m \\
        &= a_1 v_1 + a_2 (v_1+v_2) + \cdots + a_m (v_1 + \cdots + v_m) \\
        &= (a_1 + \cdots + a_m) v_1 + (a_2 + \cdots + a_m) v_2 + \cdots + a_m v_m.
    \end{aligned}
    \]
    Thus, $u \in \myspan{v_1, \ddd, v_m}$, and therefore,
    \begin{equation}
      \label{eq: fact2}
      \myspan{w_1, \ddd, w_m} \subseteq \myspan{v_1, \ddd, v_m}.
    \end{equation}

    Combining \eqref{eq: fact1} and \eqref{eq: fact2} shows that $\myspan{w_1, \ddd, w_m} = \myspan{v_1, \ddd, v_m}$.
  \end{xprf}
\end{xrcs}

\exercise{11}
\begin{xrcs}
  Suppose $v_1, \ddd, v_m$ is a linearly independent list in $V$ and $w \in V$. Prove that $v_1, \ddd, v_m, w$ is linearly independent if and only if
  \begin{equation}
    w \notin \myspan{v_1, \ddd, v_m}.
  \end{equation}
  \begin{xprf}
    \Rightarrowdirection Suppose $v_1, \dots, v_m, w$ is linearly independent. If $w \in \myspan{v_1, \dots, v_m}$, then $w = c_1 v_1 + \dots + c_m v_m$, which gives a non-trivial linear combination $c_1 v_1 + \dots + c_m v_m - w = 0$ with coefficient $-1 \neq 0$ for $w$, contradicting linear independence. Thus $w \notin \myspan{v_1, \dots, v_m}$.

    \Leftarrowdirection Conversely, suppose $w \notin \myspan{v_1, \dots, v_m}$. Consider scalars $a_1, \dots, a_m, c \in \myF$ such that
    \[
      a_1 v_1 + \dots + a_m v_m + c w = 0.
    \]
    If $c \neq 0$, then $w = -\frac{a_1}{c}v_1 - \dots - \frac{a_m}{c}v_m \in \myspan{v_1, \dots, v_m}$, a contradiction.
    Thus $c = 0$. The equation then reduces to $a_1 v_1 + \dots + a_m v_m = 0$.
    Since $v_1, \dots, v_m$ is linearly independent, $a_1 = \dots = a_m = 0$.
    Therefore, all coefficients are zero, and $v_1, \dots, v_m, w$ is linearly independent.
  \end{xprf}
\end{xrcs}

\begin{ainote}[Why the order of the two steps matters]
  In the $\Leftarrow$ direction it is essential that $c = 0$ is established \emph{first}. Only once $c$ is known to vanish does the equation collapse to $a_1 v_1 + \cdots + a_m v_m = 0$, where the hypothesis on $v_1, \ddd, v_m$ can be applied. Trying to use linear independence of $v_1, \ddd, v_m$ while the term $cw$ is still present proves nothing, since that hypothesis says nothing about $w$.

  This exercise is the engine behind the basis-extension arguments later on: it says that a vector may be appended to a linearly independent list without spoiling independence exactly when it is genuinely new, i.e., not already reachable from the vectors present.
\end{ainote}

\exercise{16}
\begin{xrcs}
  Explain why no list of four polynomials spans $\polyn_4(\myF)$.

  \begin{xsol}
    The list $1, z, z^2, z^3, z^4$ is linearly independent in $\polyn_4(\myF)$ and has length $5$. By \ref{thm: length of linearly independent list at most length of spanning list}, the length of every linearly independent list in a finite-dimensional vector space is less than or equal to the length of every spanning list.
    Therefore, every spanning list of $\polyn_4(\myF)$ must have length at least $5$. Thus, no list of four polynomials can span $\polyn_4(\myF)$.
  \end{xsol}
\end{xrcs}

\exercise{19}
\begin{xrcs}
  Suppose $p_0, p_1, \ddd, p_m$ are polynomials in $\polyn_m(\myF)$ such that $p_k(2) = 0$ for $0 \leq k \leq m$.
  Prove that $p_0, p_1, \ddd, p_m$ is not linearly independent in $\polyn_m(\myF)$.

  \begin{xprf}
    Define the evaluation functional $\varphi : \polyn_m(\myF) \to \myF$ by $\varphi(p) := p(2)$ for all $p \in \polyn_m(\myF)$.
    The polynomials that vanish at $z=2$ lie in the subspace $\mynull(\varphi)$ of $\polyn_m(\myF)$.
    Any polynomial $p \in \mynull(\varphi)$ satisfies $p(2) = 0$, so $p(z) = (z-2) q(z)$ for some $q \in \polyn_{m-1}(\myF)$.
    Conversely, every $q \in \polyn_{m-1}(\myF)$ produces a polynomial $(z-2)q(z)$ of degree at most $m$ vanishing at $2$, so the map
    \begin{equation}
      \Phi : \polyn_{m-1}(\myF) \to \mynull(\varphi), \qquad \Phi(q) := (z-2)\,q(z)
    \end{equation}
    is onto. It is also injective, since $(z-2)q(z)$ is the zero polynomial only for $q = 0$. So $\Phi$ is an isomorphism, and the two spaces have the same dimension.
    Since $\dim \polyn_{m-1}(\myF) = m$, we have $\dim \mynull(\varphi) = m$.
    The list $p_0, p_1, \ddd, p_m$ consists of $m+1$ vectors in the $m$-dimensional subspace $\mynull(\varphi)$.
    By Theorem \ref{thm: length of linearly independent list at most length of spanning list}, no linearly independent list in $\mynull(\varphi)$ can have length greater than $\dim \mynull(\varphi) = m$.
    Therefore, the list $p_0, p_1, \ddd, p_m$ of length $m+1$ is linearly dependent (not linearly independent).
  \end{xprf}

  \begin{failedattempt}
    Define the polynomials $p_0, p_1, \ddd, p_m$ for all $a$'s as elements of $\myF$:
    \[
    \begin{aligned}
      p_0(z) &:\equiv a_{0,0}  \\
      p_1(z) &:\equiv a_{1,0} + a_{1,1} z  \\
      p_2(z) &:\equiv a_{2,0} + a_{2,1} z + a_{2,2} z^2  \\
      \vdots & \quad \\
      p_m(z) &:\equiv a_{m,0} + a_{m,1} z + a_{m,2} z^2 + \cdots + a_{m,m} z^m
    \end{aligned}
    \]

    By assumption, $p_k(2) = 0$ for $k = 0, \ddd, m$, which gives the following constraints:
    \[
    \begin{aligned}
      p_0(2) &= a_{0,0} = 0 \\
      p_1(2) &= a_{1,0} + a_{1,1} 2 = 0 \\
      p_2(2) &= a_{2,0} + a_{2,1} 2 + a_{2,2} 2^2  = 0  \\
      \vdots & \quad \\
      p_m(2) &= a_{m,0} + a_{m,1} 2 + a_{m,2} 2^2 + \cdots + a_{m,m} 2^m  = 0
    \end{aligned}
    \]

    From the definition of linear independence, we examine the equation:
    \[
      b_0 p_0 + b_1 p_1 + b_2 p_2 + \cdots + b_m p_m = 0.
    \]

    Substituting $p_0(2) = 0$, this reduces to:
    \[
      b_0 \cdot 0 + b_1 p_1 + b_2 p_2 + \cdots + b_m p_m = 0.
    \]

    Any choice of $b_0$ will solve this equation and give us a non-trivial solution.
    So the list of polynomials $p_0, p_1, \ddd, p_m$ is not linearly independent.
  \end{failedattempt}

  \begin{ainote}[Where it breaks]
    Two separate things go wrong.

    \emph{The shapes of the $p_k$ are invented.} Writing $p_0(z) :\equiv a_{0,0}$ makes $p_0$ a constant, $p_1$ of degree at most $1$, and so on. The exercise says only that each $p_k$ lies in $\polyn_m(\myF)$, so \emph{every} $p_k$ may have degree up to $m$; there is no reason the index $k$ should cap the degree. (Indeed $p_0(2)=0$ together with $p_0$ constant would force $p_0 \equiv 0$, which trivialises the problem into a case that need not occur.)

    \emph{A value at a point is confused with the zero polynomial.} The equation $b_0 p_0 + \cdots + b_m p_m = 0$ is an identity between \emph{functions}: it must hold for every $z$, not just at $z = 2$. Substituting $p_0(2) = 0$ replaces the polynomial $p_0$ by the number $0$ in an equation where $p_0$ appears as a vector. From $p_0(2)=0$ one may conclude nothing about $b_0 p_0$ as an element of $\polyn_m(\myF)$ unless $p_0$ is the zero polynomial, and it need not be.

    The step \qt{any choice of $b_0$ will solve this equation} is where both errors surface at once: it is only true of the evaluated equation at $z = 2$, which is not the equation linear independence is about.

    The repair is to stop working coefficient by coefficient and count dimensions instead, as the proof above does: all $m+1$ polynomials are trapped inside the $m$-dimensional space $\mynull(\varphi)$, and $m+1$ vectors never fit independently into $m$ dimensions.
  \end{ainote}
\end{xrcs}

\clearpage
